\section{Physics versus computation \label{sec:phys_vs_comp}}

\begin{comment}
BITS FOR THIS SECTION
%===============================

WUD 4	What does the computational approach impose? Asynchronicity, possibly. Hence my trick to get round strategy. 
WUD 24: how do physics methods and computing methods define the nature of the location problem differently; what does this mean for 'translating' one to the other? Do the different foundations make them mutually exclusive?


\end{comment}

Ideas from the history of physics can be particularly fruitful for thinking through economic and agent-based problems, if only to illustrate when physical and social modelling differ incommensurably.
%I'm thinking here of using gravity as a good example of where the explanans/explanandum distinction could go further, but that at one level it does all you need it to.
---

The structure of \gls{ge} and \gls{abm} are entirely different. They are built on different foundations, and consequently have developed separate philosophies for model development. This may seem an obvious statement, but the consequences are not always so clear. The central lesson is that a direct translation from one approach to the other is not possible. Section \ref{sec:fowler} discusses Fowler's attempt at recreating the \gls{coremodel} in agent form. Fowler does an excellent job of demonstrating the point just made: it doesn't make sense to directly translate the neoclassical core model into the computational form of an agent model.

Computational modellers are really doing exactly what nineteenth century economists were - using tools they found available at the time. As Blaug points out \fixref{blaug on Walras} when discussing Walras, one might argue there is a fine line between using those tools to help gain insights into how the economic world works, and mistaking those tools for reality. \cmt{This is perhaps the point that Friedman is making. Also, loose referencing in this section (physics vs computation)}

%I think the argument that Blaug has, looking at e.g. Walras, is that somewhere there's a line crossed between using those tools to help gain insights into how the economic world works, and mistaking the tools for reality. This is perhaps the point that Friedman is making. 
Some agent-based modellers see precedents for their ideas in the pre-computer era. Vriend, for example, wonders 'was Hayek an agent-based computational economist?' Hayek, he argues, ``insisted on the need to consider a market economy as a truly decentralised system of interacting individual agents'' (\citealt{vriend_was_2002} p.2.) and in this, Vriend and others (e.g. \citealt{miller_complex_2007}) can imagine that, had Hayek only been able to access present technology, he would surely have embraced it, and complexity theory along with it.

But is this sort of counterfactual (``what if...'') of any real value? If it done carefully, perhaps so. For the present thesis, two points are worth noting. Firstly, Hayek holds a special place in the hearts of some agent modellers: this draws attention to \gls{abm}'s tunnel-vision focus on complexity. Secondly, there's a question hidden in this line of thinking - precisely how much does the knowledge and technology of a particular era carve the channels that theory-making must take?

The current theoretical structure of \gls{ge} shares the lineage of all neoclassical economics, and can be traced back to the \gls{margrev}. That revolution was built on the liberal use of ideas from physics. Mirowski, for example, argues (paraphrased by Blaug) that ``the whole of neo-classical economics ever since the \gls{margrev} has been an attempt to create an economics that emulates all the essential features of of nineteenth-century physics.'' (\citealt{blaug_economic_1997} p.284) As Blaug notes, this is rather strong; there was no conscious attempt to emulate physics. Rather, theorists, looking around for ideas, %need some detail from e.g. blaug 1997 p.284 around here. But the point is... say which ones! 
found in Newtonian physics - a very well established and rich mathematical tradition by the nineteenth century - the tantalising possibility of producing exact economic results. The \gls{margrev} was precisely about making calculus serve as a tool for explaining all of human economic behaviour through 'marginal' rates of change.

%not sure where I'm going with this, so let's try a different tack...
The aim here isn't to develop a historiography of economic theory in any great depth; the point is that, when comparing \gls{ge} and \gls{abm} approaches, the foundation stones on which they are built can help a great deal in making that comparison transparent. As regards \gls{ge}, a theme that runs through explanations of chosen methods is compromise - the analytical gains made by finding mathematically tractable ways of describing a problem.

The tendency in \gls{abm}, and in \gls{ace} in particular, has been to label all of neoclassical economics as an analytic fantasy - perhaps understandable at a time when no computers were available, but entirely unforgiveable in the twenty-first century. 
%citation?
This is a mistake, because it means too much is lost; there is no simple shortcut that can offer a way to dismiss an entire body of theory in this way. The job of this thesis is to disassemble some of the key components that make up \gls{ge} to see how they work. Only then can a direct comparison with an agent-approach have any chance of succeeding. As will be discussed in \ref{sec:fowler}, when examining a previous attempt to re-create the core model in agent form, the foundation stones of each method mean a straightforward 'conversion' is not advisable. 



